\scanpage{30}
\unsure{Item|Idem|Itou}, si on regarde les traces de ces décompositions binaires sur $S^\ast=S\amalg A'$. Donc pour tte application \uline{surjective} $X\unsure{\mathrel{\text{---}}|\longrightarrow}S^\ast$, si on regarde \unsure{les des|les les} images inverses des décompositions binaires précédentes de $S^\ast$. \unsure{\deleted{D}}Or la donnée d'une telle application, revient à celle d'une famille $(X_\sigma)_{\sigma\in S^{\unsure{\ast|+}}}$ \inserted{d'ens $\ne\varnothing$}, i.e. \deleted{à} celle des deux familles $(U_s)_{s\in S}$ et $(E_a)_{a\in A}$ d'ens non vides.

Le cas où certains des ensembles $U_s,E_a$ pourraient être vides est à examiner avec plus de soin, il donne néanmoins ce qu'on veut sss on a les conditions $(\ast)$ $(\ast\ast)$ ci-dessus, et la condition de non-parallélisme stricte sss on a \uline{ceci} (plus fort que $(\ast\ast)$)

$(\ast\ast\ast)$ Si $s\in S$ est un sommet \uline{d'ordre 2} \deleted{tel} \inserted{et} \deleted{que $U_s=\varnothing$, et} \inserted{$a\in A$ une arête incidente, on a}
\[
U_s\ne\varnothing\ \unsure{\text{ou}|\text{--}}\ E_{\unsure{c|a}}\ne\varnothing
\]
% archive sheet number: 29
